%%%%%%%%%%%%%%%%%%%%%%%%%%%%%%%%%%%%%%%%%%%%%%%%%%%%%%%%%%%%%%
%%%%%%%%%%%%%%%%% MA-0705 Análisis Real I %%%%%%%%%%%%%%%%%%%%
%%%%%%%%%%%%%%%%%%%%%%%%%%%%%%%%%%%%%%%%%%%%%%%%%%%%%%%%%%%%%%
\newpage
\subsection*{MA-0705 Análisis Real I}
\addcontentsline{toc}{subsection}{MA-0705 Análisis Real I}

\hypertarget{MA0705}{}
\begin{refsection}
\begin{table}[H]
\centering{
\begin{tabular}{|ll|}
\hline
%\multicolumn{2}{|c|}{\textbf{Nivel de virtualidad del curso:} Virtual}\\ \hline
\textbf{Año:} IV  & \textbf{Requisitos:} MA-0635 Topología General \\ 
\textbf{Ciclo:} VII  & \textbf{Correquisitos:} Ninguno \\ 
\textbf{Tipo de curso:} Teórico & \textbf{Horas presenciales por semana:}   \\ 
\textbf{Créditos:} 4 & \textbf{Horas de trabajo independiente por semana:}  \\ \hline
\end{tabular}
}
\end{table}



\subsubsection*{Descripción del curso}

Este es un curso del tercer año de la carrera Bach. y Lic. en Matemática en el que se presentan las bases del Análisis Real. En particular, aborda la Teoría de Medida e Integración de Lebesgue en $\R^n$ y marca una transici\'{o}n entre el estudio del an\'{a}lisis de espacios eucl\'{\i}deos y la teor\'{\i}a de espacios abstractos, así como el inicio de un mayor rigor en el abordaje de los temas. 

El curso inicia con una introducci\'{o}n a la teor\'{\i}a de variación acotada y la integral de Riemann-Stieltjes, para luego enfocarse en la teor\'{\i}a de la medida e integraci\'{o}n. De esta forma se
dispone de una herramienta que permite la construcci\'{o}n de espacios normados y de Hilbert de manera concreta, sentando las bases para abordar temas avanzados de an\'{a}lisis real y funcional en cursos posteriores.

Para el buen desempeño en este curso es necesario tener conocimientos previos de la topología básica de $\R^n$, así como del cálculo diferencial e integral en varias variables.



\subsubsection*{Objetivos}

\textbf{General:} Aplicar los conceptos y resultados básicos del análisis real para la resolución de problemas. \\

\textbf{Específicos:} Al finalizar el curso se espera que la persona estudiante sea capaz de:

\begin{enumerate}
\item Resolver problemas matemáticos y de disciplinas afines utilizando los principales conceptos y resultados de la teor\'{\i}a de la medida e integraci\'{o}n de Lebesgue en $\R^n$.
\item Aplicar los conceptos b\'asicos de variaci\'on de funciones, as\'i como de la integral de Riemann-Stieltjes, en el planteamiento y resoluci\'on de problemas matem\'aticos.
\item Demostrar resultados sobre propiedades de conjuntos y funciones Lebesgue-medibles, integración de Riemann-Stieltjes y de Lebesgue.
\item Relacionar los distintos modos de convergencia y los distintos tipos de integrabilidad.
\item Indagar en la literatura sobre temas que permitan complementar su formación.
\item Comunicar ideas matemáticas de manera clara y efectiva en forma oral y escrita.
\end{enumerate}

\newpage 

\subsubsection*{Contenidos}

\begin{enumerate}
\item \textbf{La integral de Riemann-Stieljes (2 semanas):} 
\begin{enumerate}
    \item Funciones de variación acotada
    \item Definición y propiedades de la integral de Riemann-Stieltjes.
\end{enumerate}

\item \textbf{Medida de Lebesgue en $\R^n$ (4 semanas):} 

\begin{enumerate}
    \item La medida exterior de Lebesgue.
    \item Conjuntos Lebesgue medibles. Ejemplo de un conjunto no medible.
    \item Caracterización de Carathéodory de la medibilidad.
    \item Propiedades de la medida de Lebesgue.
\end{enumerate}

\item \textbf{Funciones Lebesgue medibles (2 semanas):} 
\begin{enumerate}
    \item Propiedades de las funciones Lebesgue medibles. Funciones Borel medibles.
    %\item Funciones semicontinuas.
    \item Tipos de convergencia: puntual, uniforme, casi por doquier, en medida.
    \item Teorema de Egorov, Teorema de Lusin.
\end{enumerate}


\item \textbf{La integral de Lebesgue (5 semanas):} 
\begin{enumerate}
    \item Integración de funciones positivas y sus propiedades.
    \item Integración de funciones medibles.
    \item Teoremas de convergencia. Convergencia en $L^1$.
    \item Relación con las integrales de Riemann.
    \item Teoremas de Fubini y Tonelli.
\end{enumerate}

\item \textbf{Espacios $L^p$ (2 semanas):} 
\begin{enumerate}
    %\item Espacios $\ell^p$: desigualdades, convergencia, completitud.
    \item Definición y propiedades de los espacios $L^p$: normas, Desigualdad de Hölder, convergencia, densidad y separabilidad.
    \item Aproximaciones de la identidad.
    %\item Espacios $L^p$-débiles.
\end{enumerate}

\end{enumerate}

\subsubsection*{Metodología}

\noindent \textbf{Lineamientos metodológicos:} La persona docente a cargo de este curso debe respetar las siguientes pautas:
\begin{enumerate}
    \item Fomentar el desarrollo de intuición sobre los conceptos estudiados en el curso. 
    \item Promover la comunicación clara y rigurosa de argumentos e ideas matemáticas de manera oral y escrita.
    \item Cuando la naturaleza del contenido lo permita, se deben utilizar representaciones visuales que apoyen la comprensión de los temas.
    \item Fomentar una visión integral de las matemáticas y de cómo se enlazan las diferentes áreas a través del análisis real.
\end{enumerate}

\textbf{Sugerencias metodológicas:} Adicionalmente, se tienen las siguientes recomendaciones para la persona docente a cargo de este curso:
\begin{enumerate}
    \item Repasar el Axioma de Elección previo a la discusión de conjuntos no medibles.
    \item Utilizar el recurso tecnológico para reforzar distintos conceptos estudiados en el curso por medio de la visualización, cálculo y experimentación.
    \item Aprovechar momentos adecuados del curso para motivar el estudio por medio de reseñas acerca del desarrollo histórico del análisis real y su importancia en aplicaciones dentro y fuera de la matemática. Por ejemplo se puede consultar \cite{Fol16} y \cite{Bre08}.
    \item En la evaluación, se sugiere utilizar estrategias que promuevan el trabajo constante de parte de las personas estudiantes a lo largo del ciclo lectivo. Por ejemplo, esto se puede hacer por medio de listas de ejercicios recomendados o proyectos.
    \item Para fomentar el desarrollo de las habilidades investigativas en el estudiantado, se sugiere la implementación de proyectos de investigación y participación en coloquios o seminarios de investigación.
    \item Dado que hoy en día la mayoría del trabajo en matemática, tanto a nivel académico como profesional, es de naturaleza colaborativa, se sugiere a la persona docente implementar estrategias que promuevan el trabajo en equipo.
\end{enumerate}



\nocite{Apo76}
\nocite{Lie01}
\nocite{RF10}
\nocite{Roy88}
\nocite{WZ15}
\nocite{SS05}
\nocite{Hei19}
\nocite{BBT08}
\nocite{Bre08}
\nocite{Fol99}




\printbibliography[heading=subbibliography,section=\the\value{refsection}]
\end{refsection}
